% Ad Familiares 6.21 (ad-familiares-06-21)
% Source: https://raw.githubusercontent.com/PerseusDL/canonical-latinLit/master/data/phi0474/phi056/phi0474.phi056.perseus-lat1.xml
% Fetched by scripts/fetch_latin.py

% Dateline (Perseus): Scr. Romae circ. m. Apr. a. 709 (45) .

\ciceroLetterOpener{CICERO TORANIO.}

\ciceroSection{1} etsi cum haec ad te scribebam, aut appropinquare exitus huius calamitosissimi belli aut iam aliquid actum et confectum videbatur, tamen cotidie commemorabam te unum in tanto exercitu mihi fuisse adsensorem et me tibi, solosque nos vidisse quantum esset in eo bello mali, in quo spe pacis exclusa ipsa victoria futura esset acerbissima, quae aut interitum adlatura esset, si victus esses, aut si vicisses, servitutem. itaque ego, quem tum fortes illi viri et sapientes, Domitii et Lentuli, timidum esse dicebant (eram plane; timebam enim ne evenirent ea quae acciderunt), idem nunc nihil timeo et ad omnem eventum paratus sum. Cum aliquid videbatur caveri posse, tum id neglegi dolebam; nunc vero eversis omnibus rebus, cum consilio profici nihil possit, una ratio videtur, quicquid evenerit, ferre moderate, praesertim cum omnium rerum mors sit extremum et mihi sim conscius me, quoad licuerit, dignitati rei p. consuluisse et hac amissa salutem retinere voluisse.

\ciceroSection{2} haec scripsi, non ut de me ipse dicerem, sed ut tu, qui coniunctissima fuisti mecum et sententia et voluntate, eadem cogitares. Magna enim consolatio est cum recordare, etiamsi secus acciderit, te tamen recte vereque sensisse. atque utinam liceat aliquando aliquo rei p. statu nos frui inter nosque conferre sollicitudines nostras, quas pertulimus tum cum timidi putabamur, quia dicebamus ea futura quae facta sunt.

\ciceroSection{3} de tuis rebus nihil esse quod timeas praeter universae rei p. interitum tibi confirmo; de me autem sic velim iudices, quantum ego possim, me tibi, saluti tuae liberisque tuis summo cum studio praesto semper futurum. vale .
